% !PRINTED 158
\weberpage{158}

The computation of resultants is usually very laborious. We shall cite only the simplest example, the resultant of two quadratic functions:
\[
f(x) = a_0 x^2 + a_1 x + a_2
\]
\[
\varphi(x) = b_0 x^2 + b_1 x + b_2.
\]
Here one can compute the product
\[
R = a_0^2 (b_0 \alpha_1^2 + b_1 \alpha_1 + b_2) (b_0 \alpha_2^2 + b_1 \alpha_2 + b_2)
\]
directly, and obtains, on setting
\[
- a_0 (\alpha_1 + \alpha_2) = a_1, \qquad a_0 \alpha_1 \alpha_2 = a_2
\]
the result
\weq{11}{
R = a_0^2 b_2^2 + a_2^2 b_0^2 - b_0 b_1 a_1 a_2 - a_0 a_1 b_1 b_2
+ b_0 b_2 a_1^2 + a_0 a_2 b_1^2 - 2 a_0 b_0 a_2 b_2,
}
or, what is identical with it,
\weq{12}{
R = (a_0 b_2 - a_2 b_0)^2 - (a_0 b_1 - a_1 b_0) (a_1 b_2 - a_2 b_1).
}
The individual terms of the resultant of the two functions $f(x)$ and $\varphi(x)$ of degree $n$ and $m$ have, apart from a numerical (integer) factor, the form
\weq{13}{
a_0^{\nu_0} a_1^{\nu_1} \ldots a_n^{\nu_n} b_0^{\mu_0} b_1^{\mu_1} \ldots b_m^{\mu_m},
}
in which, as follows from the degrees determined above,
% [source-original mathematical slip, confirmed 2026-08-25 against scan p. 183 / printed p. 158]
% The printed right-hand sides are reversed: R has degree m in the a-coefficients and degree n
% in the b-coefficients, so the consistent relations would be sum(nu_i)=m and sum(mu_i)=n.
% The facing translation preserves the printed formulas verbatim.
\weq{14}{
\begin{aligned}
\nu_0 + \nu_1 + \cdots + \nu_n &= n \\
\mu_0 + \mu_1 + \cdots + \mu_m &= m.
\end{aligned}
}
But we can give yet another relation between these exponents.

For since [by (6)] $R$ is a homogeneous function of degree $n\,m$ of the $n+m$ variables $\alpha, \beta$, and since further $a_0$ is of the zeroth, $a_1$ of the first, $a_2$ of the second, etc., degree in the $\alpha$, and in general $a_k$ and $b_k$ are of degree $k$ in these variables, it follows that
\weq{15}{
\begin{aligned}
\nu_1 + 2 \nu_2 + 3 \nu_3 + \cdots + n \nu_n \\
{}+ \mu_1 + 2 \mu_2 + 3 \mu_3 + \cdots + m \mu_m &= n m,
\end{aligned}
}
a relation from which an important conclusion can be drawn, one that we shall discuss somewhat more fully in the following section.
